% ===== 原文 pp.47–48（§5.4 问题集 IV）=====
% [原文 p.47]

\subsection{问题集 IV}

本问题集旨在让读者熟悉多边形与超多边形的空间，以及它们在 Higgs 丛意义下的解释。此外，读者会发现这些问题有助于理解李群之间的各种对应，以及它们所诱导的 Hitchin 系统之间的对应。

\begin{prob}{16} \textbf{同源。}为理解一个 $\Sl(4,\R)$-Higgs 丛与自身的张量积，先从一个 $\Sl(4,\R)$-Higgs 丛 $(E,\varphi)$ 出发，其中 $E$ 的秩为 4、行列式平凡，且 $\Phi$ 关于局部标架 $\{e_1,\dots,e_4\}$ 是无迹的、对称的。则 $E\otimes E$ 上的 Higgs 场为
\[
\Phi = \varphi\otimes I + I\otimes\varphi.
\]
\begin{enumerate}[label=(\roman*),leftmargin=2.2em]
\item 回忆：对任意秩为 $n$ 的向量丛 $\mathbf{V}$ 以及任意 $k<n$，都有同构
\[
\Lambda^{k}(\mathbf{V}^{*})\otimes\Lambda^{n}(\mathbf{V})\longrightarrow\Lambda^{n-k}(\mathbf{V}).
\]
当取 $n=4$、$k=2$、行列式 $\Lambda^{4}\mathbf{V}\simeq\cO$，并且 $\mathbf{V}$ 具有由同构 $q:\mathbf{V}\simeq \mathbf{V}^{*}$ 给出的正交结构时，会出现什么结构？
% [Tutor Note: p.47 问题 16.i 的同构 $\Lambda^{k}(V^{*})\otimes\Lambda^{n}(V)\to\Lambda^{n-k}(V)$ 由「用体积形式做缩并」给出：$\Lambda^{k}V^{*}$ 与 $\Lambda^{n}V$ 配对后落在 $\Lambda^{n-k}V$ 上。取 $n=4$、$k=2$ 且 $\Lambda^{4}V\simeq\cO$，该同构就成为 $\Lambda^{2}(V^{*})\simeq\Lambda^{2}(V)$；再配合 $q:V\simeq V^{*}$，$\Lambda^{2}V$ 上便诱导出正交结构（自对偶/反自对偶分解），这正是下一小问中 $\Lambda^{2}E$ 上 $G$-Higgs 场结构的来源。]
\item 关于
\[
E\otimes E = \Lambda^{2}E\oplus\Sym^{2}(E),
\]
的局部标架，$\Phi$ 到 $\Lambda^{2}E$ 上的限制是什么样子？它是否具有某个实群 $G$ 的 $G$-Higgs 场的结构？
\item 计算整个 Higgs 场 $\Phi$ 在问题 16.ii. 的标架下是什么样子，亦即关于
\[
\{e_i\wedge e_j\mid i<j\}\cup\{e_i\otimes e_i\mid i=1,2,3,4\}\cup\{e_i\otimes e_j+e_j\otimes e_i\mid i<j\}
\]
是什么样子。
\end{enumerate}
% [Tutor Note: p.47 问题 16 首段先以 $(E,\varphi)$ 引入 Higgs 丛、随即写 "$\Phi$ is trace free"，同一段末又令 $\Phi=\varphi\otimes I+I\otimes\varphi$（$E\otimes E$ 上的新 Higgs 场）。按记号 $\varphi$ 与 $\Phi$ 的分工（$\varphi$ 在 $E$ 上、$\Phi$ 在 $E\otimes E$ 上），"$\Phi$ is trace free" 中的 $\Phi$ 应为 $\varphi$；又首句 "an $\Sl(4,\R)$-Higgs bundles" 的 "bundles" 亦为笔误（应为单数）。正文照原样译出，仅在此指出。]
\end{prob}

\begin{prob}{17} \textbf{同源。}若 (5.2) 中的同态 $\Psi$ 限制为实形式之间的映射，且保持李代数的 Cartan 分解，则它诱导实 Higgs 丛之间的映射。考虑复群之间的态射，计算以下内容：
\begin{enumerate}[label=(\roman*),leftmargin=2.2em]
\item 说明 $\Sl(4,\C)$ 与 $\SO(6,\C)$ 之间的同源，在实辛 Higgs 丛上考虑时意味着什么；这里的实辛 Higgs 丛即丛为 $V\oplus V^{*}$、且 Higgs 场的元素对称的那些对 $(E,\varphi)$，此时
\[
\varphi = \begin{pmatrix} 0 & \beta \\ \gamma & 0 \end{pmatrix}.
\]
% [Tutor Note: p.47 问题 17.i 中 "this over pairs" 语意不通，应为 "i.e., over pairs"（或 "this is, over pairs"）之误，即「亦即丛为 $V\oplus V^{*}$ 的那些对」；译文按此理解处理。$\varphi$ 的矩阵写法 $\begin{pmatrix}0&\beta\\ \gamma&0\end{pmatrix}$ 与 $V\oplus V^{*}$ 的分解相容，$\beta,\gamma$ 分别为 $\Hom(V^{*},V)$、$\Hom(V,V^{*})$ 中的对称元素。]
\item 通过复李群之间的同源得到的 $\Sl(4,\R)$-Higgs 丛与 $\SO(3,3)$-Higgs 丛，其特征值之间有什么关系？
% [Tutor Note: p.47 问题 17.ii 的 "obtained thought the isogeny" 中 "thought" 系 "through" 之误（同音误写）；正文译为「通过……得到的」。]
\end{enumerate}
\end{prob}

% [原文 p.48]
\begin{prob}{18} \textbf{多边形。}称 $\PP^{1}$ 上的 $n$ 个有序点 $(p_1,\dots,p_n)$ 是 $r$-半稳定的，如果对任意 $p_0\in\PP^{1}$ 都有
\[
\sum_{i\,\mid\,p_i=p_0} r_i \;\leqslant\; \sum_{i\,\mid\,p_i\neq p_0} r_i .
\]
\begin{enumerate}[label=(\roman*),leftmargin=2.2em]
\item 证明：$\PP^{1}$ 上 $r$-半稳定的 $n$ 个有序点、模去射影直线的自同构所得的模空间，可以用 GIT 构造为 $\Sl_2$ 在 $(\PP^{1})^{n}$ 上作用的商，所取的是与这些权重相应的丰沛线性化（例见 [66, 第 5.2 节]）。
\item 对固定的星形箭图以及给定的 $d,n$ 值，描述与多边形空间相对应的表示空间。
\item 证明星形箭图 $Q_n$ 的每个自同构都必须固定 $v_0$，于是 $\mathrm{textrmAut}(Q_n) = \operatorname{Aut}^{+}(Q_n) \cong S_n$。
% [Tutor Note: p.48 问题 18.iii 原文 LaTeX 误写：源码中写作 textrmAut（漏掉反斜杠与花括号），PDF 上直接印出罗马体的 "textrmAut(Q_n)"。正确的排版应为 $\operatorname{Aut}(Q_n)$。正文按原文（印出的字样）照排，此处说明其本意。另注意同行的 $\operatorname{Aut}^{+}(Q_n)$ 才是作者想与 $\operatorname{Aut}(Q_n)$ 相比较的对象。]
\item 构造 $\operatorname{Aut}(Q_n)$ 的两个子群，描述它们在 $Q_n$ 上的作用，以及在相应群作用下的商箭图。
\end{enumerate}
% [Tutor Note: p.48 问题 18 的 $r$-半稳定定义可等价地写成：对每个 $p_0$，落在 $p_0$ 处的权重之和不超过总权重 $\sum_i r_i$ 的一半。这一形式与 GIT 中权重的「重心」判据一致，可直接用于核对本问题 i 中「与权重相应的丰沛线性化」。]
\end{prob}

\begin{prob}{19} \textbf{超多边形。}考虑箭图 $Q$：它有两个顶点 $1$ 与 $2$，以及两条箭头，其中 $x$ 从顶点 $1$ 指向顶点 $2$，$y$ 从顶点 $2$ 指向顶点 $1$（原文此处以上下两条相对的弧线示意该箭图，$x$ 在上、$y$ 在下；本译本按排版约定以文字描述代替图示）。
\begin{enumerate}[label=(\roman*),leftmargin=2.2em]
\item 设 $\sigma$ 是把顶点 $1$ 送到顶点 $2$、把箭头 $x$ 送到箭头 $y$ 的协变对合。何时 $d=(d_1,d_2)$ 是 [66] 意义下的 $\sigma$-相容维数向量？
% [Tutor Note: p.48 问题 19.i 原文 "the covariant involution which sending the vertex 1 to the vertex 2" 中 "which sending" 应为 "which sends"（或 "sending"），属语法笔误；译文按其本意译出。]
\item 对某个特定的复或实李群 $G_{\C}$，把 $\PP^{1}$ 上的 $G_{\C}$-Higgs 丛 $(E,\Phi)$ 的信息编码到上述箭图之中。
\item 用超多边形的语言，描述与例 5.8 相联系的抛物 Higgs 丛的模空间及其超凯勒结构？
\end{enumerate}
\end{prob}

\begin{prob}{20} \textbf{超多边形。}在理解了多边形与超多边形如何用 Higgs 丛来表述之后，考虑带框的 Jordan 箭图（framed Jordan quiver）$\mathcal{J}$：它有两个顶点 $1$ 与 $\mathrm{l}(1)$，顶点 $1$ 上有两条回路 $a$（上）与 $b$（下），另有箭头 $x:1\to\mathrm{l}(1)$ 与 $y:\mathrm{l}(1)\to 1$（原文同样以图示给出，本译本按排版约定以文字描述代替）。
\begin{enumerate}[label=(\roman*),leftmargin=2.2em]
\item 仿照 [42]，给出与这类箭图相联系的实矩映射与复矩映射（相差一个整体常数倍）的表达式。
\item 设 $\sigma:\mathcal{J}\to\mathcal{J}$ 是 Jordan 箭图的对合，它交换两个顶点、并交换相应的箭头。通过它在模空间上的诱导作用，描述不动点集。
\item 如前所述，前几节所研究的超多边形空间带有超凯勒结构，因此可以考虑三个固定的复结构及与之相伴的辛形式，并考虑它们与给定子空间的关系。仿照 [66] 与 [85]，判断问题 19.ii. 中的不动点集是否为膜；若是，给出其类型。
\end{enumerate}
% [Tutor Note: p.48 问题 20 的箭图右侧顶点在原图中标注为罗马体的 $\mathrm{l}(1)$（小写字母 l 加括号），并非数字 1；绘制时顶点 1 位于左侧、该顶点位于右侧，回路 $a$、$b$ 均以顶点 1 为基点，箭头 $y$ 指向顶点 1、$x$ 指向 $\mathrm{l}(1)$。按本译本「不重画箭图」的约定，此处以文字逐项描述，读者可对照原文 p.48 的图。]
\end{prob}
